\documentclass[11pt]{article}
\usepackage[margin=1in]{geometry}
\usepackage{amsmath,amssymb}
\usepackage{enumitem}
\usepackage{hyperref}
\usepackage{graphicx}
\usepackage{booktabs}
\usepackage{parskip}
\usepackage[T1]{fontenc}
\usepackage[utf8]{inputenc}
\usepackage{lmodern}

\title{Formal RL Length Generalization: Model Architecture, Objectives, and Training Strategy}
\author{Project Notes}
\date{June 2026}

\begin{document}

\maketitle

\section{Overview}
This document summarizes the core ideas behind the project:\par
\emph{formal\_rl\_length\_generalization}. The goal is to study whether small autoregressive Transformer policies can learn algorithmic behavior for formal-language recognition and generalize to longer input lengths than those seen during training.

The project uses symbolic traces and oracle rewards for tasks such as:
\begin{itemize}[leftmargin=*]
  \item parity
  \item modular counting with $k=3$ or $k=5$
  \item $a^*b^*$
  \item $a^n b^n$
\end{itemize}

Training lengths are restricted to $n \in \{1,2,\dots,40\}$, while evaluation also includes out-of-distribution buckets $41$--$80$, $81$--$160$, $161$--$320$, and $321$--$500$.

\section{Problem Setup}
For each example, the model receives a prompt of the form
\[
\text{TASK } \langle task\_name \rangle \text{ INPUT } x_1 x_2 \cdots x_n \text{ TRACE}
\]
It must generate a symbolic trace and final answer:
\[
s_1 s_2 \cdots s_n \text{ FINAL } y_{final}.
\]

The trace is an explicit intermediate computation, not just the final label.

\section{Reward Decomposition}
The RL reward is decomposed into two parts:
\begin{align*}
R_{process} &= \text{fraction of intermediate trace states that are correct},\\
R_{terminal} &= \mathbf{1}[\text{final answer is correct}].
\end{align*}

The combined reward used by the project is
\[
R = w_p R_{process} + w_t R_{terminal},
\]
where $w_p$ is the process weight and $w_t$ is the terminal weight.

The three reward modes are:
\begin{align*}
R &= R_{terminal} \qquad \text{(terminal only)},\\
R &= R_{process} \qquad \text{(process only)},\\
R &= w_p R_{process} + w_t R_{terminal} \qquad \text{(process + terminal)}.
\end{align*}

\section{Task-Specific Reward Forms}

\subsection{Parity}
The hidden parity state is
\[
p_0 = 0, \qquad p_t = p_{t-1} \oplus x_t.
\]
The process reward is
\[
R_{process} = \frac1n \sum_{t=1}^{n} \mathbf{1}[\hat p_t = p_t].
\]
The terminal reward is
\[
R_{terminal} = \mathbf{1}[\hat p_n = p_n].
\]

\subsection{Modular Counting}
The state is
\[
c_0 = 0, \qquad c_t = (c_{t-1}+x_t) \bmod k.
\]
The final answer is ACCEPT iff $c_n=0$.
The process reward is
\[
R_{process} = \frac1n \sum_{t=1}^{n} \mathbf{1}[\hat c_t = c_t].
\]

\subsection{a$^*$b$^*$}
The DFA states are
\[
S_0,\ S_1,\ S_2,
\]
with transitions
\begin{align*}
\delta(S_0,a)&=S_0,\quad \delta(S_0,b)=S_1,\\
\delta(S_1,b)&=S_1,\quad \delta(S_1,a)=S_2,\\
\delta(S_2,a)&=S_2,\quad \delta(S_2,b)=S_2.
\end{align*}
The process reward is the fraction of correctly predicted DFA states along the trace.

\subsection{$a^n b^n$}
The implementation tracks
\[
phase_t \in \{A\_PHASE,	ext{ }B\_PHASE,	ext{ }DEAD\},
\qquad balance_t = count_a_t - count_b_t.
\]
A per-step reward is
\[
 r_t = 0.5 R_{phase}(t) + 0.5 R_{balance}(t),
\]
where
\[
R_{phase}(t)=\mathbf{1}[\hat phase_t = phase_t],
\]
and
\[
R_{balance}(t)=1-\frac{\min(|\hat balance_t-balance_t|,cap)}{cap},\qquad cap=500.
\]
Then
\[
R_{process} = \frac1n \sum_{t=1}^{n} r_t.
\]

\section{Model Architecture}
The policy is a causal Transformer.

For input tokens $z_1,\dots,z_T$ the model forms
\[
 h_t^{(0)} = E_{token}(z_t) + E_{pos}(t).
\]
The causal attention mask ensures that position $t$ can only attend to positions $\le t$.

After the Transformer stack, the final hidden state is $h_t$.

The policy head outputs logits
\[
logits_t = W_\pi h_t + b_\pi,
\]
so that
\[
\pi_\theta(a_t \mid z_{\le t}) = \mathrm{softmax}(logits_t).
\]

The value head produces
\[
V_\theta(z_{\le t}) = W_V h_t + b_V.
\]
The value head is used in PPO; the policy head is used in SFT, PPO, and GRPO.

\section{Autoregressive Generation}
Given a prompt $q$, the model samples
\[
 a_t \sim \pi_\theta(\cdot \mid q,a_1,\dots,a_{t-1}).
\]
The full completion probability is
\[
\pi_\theta(a_{1:m}\mid q)=\prod_{t=1}^{m}\pi_\theta(a_t\mid q,a_{<t}).
\]
The log-probability is
\[
\log \pi_\theta(a_{1:m}\mid q)=\sum_{t=1}^{m}\log \pi_\theta(a_t\mid q,a_{<t}).
\]
Generation stops when $\langle eos \rangle$ is generated or when the token budget is exhausted.

\section{Supervised CoT Baseline (SFT)}
The SFT objective is cross-entropy over the teacher-forced oracle trace and final answer:
\[
L_{SFT}(\theta)= -\frac1M\sum_{t=1}^{M}\log \pi_\theta(y_t \mid q, y_{<t}).
\]
The token accuracy is
\[
Acc_{token} = \frac1M \sum_{t=1}^{M}\mathbf{1}[\arg\max_a \pi_\theta(a\mid q,y_{<t})=y_t].
\]

\section{RL Objective}
The RL objective is
\[
J(\theta)=\mathbb E_{\tau\sim \pi_\theta}[R(\tau)].
\]
The policy-gradient identity is
\[
\nabla_\theta J(\theta)=\mathbb E\Big[\sum_t \nabla_\theta \log \pi_\theta(a_t\mid s_t) A_t\Big].
\]

\subsection{PPO}
The probability ratio is
\[
\rho_t(\theta)=\frac{\pi_\theta(a_t\mid s_t)}{\pi_{\theta_{old}}(a_t\mid s_t)}.
\]
The clipped PPO objective is
\[
L_{clip}(\theta)=\mathbb E\Big[\min\big(\rho_t A_t,	ext{ clip}(\rho_t,1-\epsilon,1+\epsilon)A_t\big)\Big].
\]
The policy loss is
\[
L_{policy}=-L_{clip}.
\]
The value loss is
\[
L_{value}=\mathbb E[(V_\theta(s_t)-R_t)^2].
\]
The entropy bonus is
\[
H=-\sum_a \pi_\theta(a\mid s_t)\log \pi_\theta(a\mid s_t).
\]
The full PPO loss is
\[
L_{PPO}=L_{policy}+c_v L_{value}-c_H H.
\]

\subsection{GRPO}
For one prompt, GRPO samples multiple completions $a_{i,1},\dots,a_{i,K}$ and computes rewards $R_{i,k}$.
The group mean is
\[
\mu_i = \frac1K \sum_{k=1}^{K} R_{i,k},
\]
and the group standard deviation is
\[
\sigma_i = \sqrt{\frac1K \sum_{k=1}^{K}(R_{i,k}-\mu_i)^2 }.
\]
The group-relative advantage is
\[
A_{i,k}=\frac{R_{i,k}-\mu_i}{\sigma_i+\epsilon}.
\]
The GRPO policy loss is
\[
L_{GRPO}^{policy}= -\mathbb E\Big[\min(\rho_{i,k,t}A_{i,k},\text{clip}(\rho_{i,k,t},1-\epsilon,1+\epsilon)A_{i,k})\Big],
\]
with the total loss
\[
L_{GRPO}=L_{GRPO}^{policy}-c_H H.
\]

\section{Training Strategy}
The training pipeline has four main stages:
\begin{enumerate}[leftmargin=*]
  \item sample random task instances of length $1\dots 40$,
  \item generate a completion with the current policy,
  \item compute oracle reward from the formal task,
  \item update the policy with SFT, PPO, or GRPO.
\end{enumerate}

Key practical factors are:
\begin{itemize}[leftmargin=*]
  \item reward weighting $w_p$ and $w_t$,
  \item clip epsilon $\epsilon$,
  \item entropy coefficient $c_H$,
  \item value coefficient $c_v$,
  \item temperature and max token budget,
  \item number of PPO epochs or GRPO samples per prompt.
\end{itemize}

\section{Evaluation}
The evaluation reports three main quantities:
\begin{align*}
\text{Process score} &= \text{mean } R_{process},\\
\text{Terminal score} &= \text{mean } R_{terminal},\\
\text{Exact score} &= \mathbf{1}[\text{generated trace + final answer exactly matches target}].
\end{align*}

The evaluation buckets are:
\begin{itemize}[leftmargin=*]
  \item $1$--$40$ (training lengths),
  \item $41$--$80$,
  \item $81$--$160$,
  \item $161$--$320$,
  \item $321$--$500$.
\end{itemize}

\section{Why the Project Matters}
The central hypothesis is that dense, intermediate process rewards help the model learn the true transition rule
\[
 s_t = f(s_{t-1},x_t),
\]
which should improve length generalization beyond the train lengths.

If the model only learns a shortcut for final correctness, it may perform well on $1$--$40$ but fail on longer inputs. Process rewards are intended to make the learned algorithm more faithful to the underlying formal computation.

\section{Important Limitations}
The current codebase is a research scaffold rather than a full production RLHF implementation. Important limitations include:
\begin{itemize}[leftmargin=*]
  \item rewards are broadcast at the sequence level,
  \item process rewards are computed after generation rather than as token-aligned rewards,
  \item the PPO implementation does not yet use full reward-to-go / GAE,
  \item there is no reference-model KL penalty,
  \item strict formatting can make exact-match evaluation fragile.
\end{itemize}

\end{document}
